% TOP_PROBLEM: {"id":30007680,"problem_number":"LOCAL-30007680","title":"Algorithmic pricing conduct","category_id":21,"category":{"id":21,"name":"economics","display_name":"Economics","description":"Open economic theory statements, published research agendas, and proposed scoped specifications; question provenance and historical priority are qualified per record.","slug":"economics","order_index":21,"created_at":"2026-10-08T00:00:00Z"},"status":"research_question","external_url":null,"legacy_ids":[],"published":false,"statement":"Find conditions under which decentralized learning programs sustain materially elevated prices despite algorithm replacement, demand changes and entry.","economics":{"original_id":169,"field":"Industrial organization and competition","kind":"theoretical","formalization_basis":"adapted_research_question","source_status":"explicit_open","priority_status":"earliest_found","priority_note":"Earliest relevant explicit question or agenda passage located in this audit. No claim of first-ever formulation; earlier versions and later solutions were not exhaustively traced.","earliest_verified_reference":{"authors":["Carl Shapiro","Ali Yurukoglu"],"title":"Trends in Competition in the United States: What Does the Evidence Show?","year":2024,"url":"https://web.stanford.edu/~ayurukog/trendsincompetition.pdf","doi":"10.3386/w32762","locator":"Section 6, printed p. 35, paragraph on sophisticated pricing algorithms","evidence_summary":"Asks whether sophisticated pricing raises or lowers consumer and total surplus."},"current_reference":{"authors":["Carl Shapiro","Ali Yurukoglu"],"title":"Trends in Competition in the United States: What Does the Evidence Show?","year":2024,"url":"https://web.stanford.edu/~ayurukog/trendsincompetition.pdf","doi":"10.3386/w32762","locator":"Section 6, printed p. 35, paragraph on sophisticated pricing algorithms","evidence_summary":"Asks whether sophisticated pricing raises or lowers consumer and total surplus."},"formalization_note":"The source verifies a broader question or research agenda; the population, estimand, equations and restrictions here are our scoped adaptation. This is not a quoted canonical conjecture, and the audit does not prove contemporary unsolved status for every special case. Independent math review specifies the horizon family, incumbent/entrant algorithm roles, and benchmark selection."},"statusQualification":"Published question or research agenda verified at the cited locator. The mathematical specification is adapted for this catalogue and includes additional assumptions; source publication does not establish first-ever priority or certify this exact model remains unsolved. The source verifies a broader question or research agenda; the population, estimand, equations and restrictions here are our scoped adaptation. This is not a quoted canonical conjecture, and the audit does not prove contemporary unsolved status for every special case. Independent math review specifies the horizon family, incumbent/entrant algorithm roles, and benchmark selection.","statusReviewedAt":"2026-10-08"}
% FIRST_POSTED: 2026-10-08
% ENTRY_KIND: statement_only
% !TeX program = lualatex
\documentclass[11pt]{article}
\usepackage{fontspec}
\usepackage[a4paper,margin=25mm]{geometry}
\usepackage{amsmath,amssymb,mathtools}
\usepackage{xurl}
\usepackage[hidelinks]{hyperref}
\newcommand{\sref}[2]{\hyperlink{source:#1:#2}{[S#2]}}
\setlength{\emergencystretch}{3em}
\begin{document}
% problemId:30007680
\section*{Algorithmic pricing conduct}\label{30007680}
\subsection{Definitions and mathematical statement}
Consider a finite-horizon market with \(N\geq2\) incumbent firms and periods \(t=1,\ldots,T\), evaluated over increasing positive integer horizons \(T\). Each active firm sets a price in a common compact interval, faces a specified demand function, and has observed production costs. Pricing programs learn from past public prices, quantities and their own profits; they exchange no direct messages. Potential entrants face a specified entry-cost distribution and decide when to enter. For an algorithm class \(\mathcal A\), define \(\bar p(a)\) as the quantity-weighted mean price generated by algorithm profile \(a\in\mathcal A^N\), averaged over learning randomness and demand shocks. Let \(p^B\) be the corresponding price benchmark from the same environment with independently optimal one-period pricing and identical entry opportunities and a specified equilibrium-selection rule; entrant algorithms are fixed separately in both comparisons. For a prespecified economically material threshold \(\varepsilon>0\), identify conditions under which \(\bar p(a)-p^B\geq\varepsilon\) persists as \(T\) increases and survives unilateral algorithm replacement. Include perturbations of demand, exploration rates and entrant behavior, and distinguish reciprocal punishment from unilateral pricing commitment. Report which conditions depend on algorithm architecture or market stationarity. This proposed dynamic specification extends a research agenda about sophisticated pricing; it is not a claim that a particular learning program necessarily colludes or that elevated prices prove communication.

\par\textbf{Formulation and scope.} The source verifies a broader question or research agenda; the population, estimand, equations and restrictions here are our scoped adaptation. This is not a quoted canonical conjecture, and the audit does not prove contemporary unsolved status for every special case. Independent math review specifies the horizon family, incumbent/entrant algorithm roles, and benchmark selection.

\par\textbf{Open-question provenance.} An explicit question or research agenda was verified at the source locator. This does not by itself certify that every later version remains unresolved \sref{30007680}{1}.

\par\textbf{Historical priority.} Earliest relevant explicit question or agenda passage located in this audit. No claim of first-ever formulation; earlier versions and later solutions were not exhaustively traced.

\subsection{Short English statement}
Find conditions under which decentralized learning programs sustain materially elevated prices despite algorithm replacement, demand changes and entry.

\subsection{Sources}
\begin{itemize}\raggedright
\item[\textnormal{[S1]}]\hypertarget{source:30007680:1}{} Carl Shapiro, Ali Yurukoglu. 2024. Trends in Competition in the United States: What Does the Evidence Show?. Section 6, printed p. 35, paragraph on sophisticated pricing algorithms.
\url{https://web.stanford.edu/~ayurukog/trendsincompetition.pdf}.
DOI: 10.3386/w32762.
{\small\textit{Earliest verified question/agenda reference; historical priority qualified below. Asks whether sophisticated pricing raises or lowers consumer and total surplus.}}
\end{itemize}
\end{document}
